\chapter{Допуск нелокального шеститочечного барионного ядра}
\label{ch:v8-baryon-nonlocal-six-point-kernel-admission-gate}

Локальные беспроизводная и производная ветви статического $W_3$ закрыты.
Это не запрещает нелокальный интегральный оператор. Отделим уже выведенный
внутренний цветовой носитель $W_3$ от импульсной формы и проверим минимальный
контракт без подбора наблюдаемой массы.

Пусть $z=k^2\geq0$ --- евклидов квадрат переданного импульса. Для
вещественного форм-фактора $f$ определим
\begin{equation}
 K_f(k)=\lambda_3 f(k^2)W_3.
 \label{eq:v8-baryon-nonlocal-kernel-definition}
\end{equation}
Самосопряжённость, калибровочная и перестановочная инвариантность, а также
нулевые частичные следы наследуются от $W_3$ при вещественном $f$.

Два точных положительных кандидата имеют вид
\begin{equation}
 f_1(z)=\frac1{z+1},
 \qquad
 f_2(z)=\frac1{2(z+1)}+\frac2{z+4}.
 \label{eq:v8-baryon-nonlocal-two-form-factors}
\end{equation}
Они одинаково нормированы на статическом срезе, но различны вне его:
\begin{equation}
 f_1(0)=f_2(0)=1,
 \qquad
 f_2(z)-f_1(z)=\frac{3z}{2(z+1)(z+4)}>0\quad(z>0).
 \label{eq:v8-baryon-nonlocal-static-normalization}
\end{equation}
В частности,
\begin{equation}
 f_1'(0)=-1,
 \qquad f_2'(0)=-\frac58,
 \qquad f_1(1)=\frac12,
 \qquad f_2(1)=\frac{13}{20}.
 \label{eq:v8-baryon-nonlocal-slope-witness}
\end{equation}

Неединственность непрерывна: для $0\leq t\leq1$
\begin{equation}
 f_t(z)=(1-t)f_1(z)+t f_2(z)
 \label{eq:v8-baryon-nonlocal-convex-family}
\end{equation}
положителен при $z\geq0$ и удовлетворяет $f_t(0)=1$.
Оба представителя имеют положительное спектральное разложение
\begin{equation}
 f(z)=\sum_j\frac{r_j}{z+m_j^2},
 \qquad
 (r_j,m_j^2)=(1,1)
 \quad\hbox{или}\quad
 \left(\frac12,1\right),(2,4).
 \label{eq:v8-baryon-nonlocal-positive-spectral-data}
\end{equation}
Таким образом, они допускают обычное чтение как эффективные обменные
каналы с положительными квадратами масс и вычетами. Но ни массы, ни вычеты,
ни число полюсов не содержатся в прежнем $42$-мерном конечном родителе.

Нелокальность здесь проверяема алгебраически: символ локального
конечнопорядкового дифференциального ядра полиномиален по $z$, тогда как
\begin{equation}
 \operatorname{poles}(f_1)=\{-1\},
 \qquad
 \operatorname{poles}(f_2)=\{-1,-4\}.
 \label{eq:v8-baryon-nonlocal-pole-test}
\end{equation}
Следовательно, $K_f$ является обратнокинетическим, то есть эффективным
нелокальным ядром, а не скрытым локальным кубическим членом.

На трёхчастичной амплитуде допустимый операторный уровень теперь записан как
\begin{equation}
 \Psi(p)=\int K_f(p-q)\Psi(q)\,dq
 +\sum_{a=1}^{3}\int K_{(2)}^{(a)}(p,q)\Psi(q)\,dq.
 \label{eq:v8-baryon-nonlocal-faddeev-placement}
\end{equation}
Это размещает $W_3$ в неприводимой трёхчастичной части, но ещё не доказывает
полюс, нормировку амплитуды или наблюдаемую массу.

\begin{theorem}[Математический допуск нелокального класса]
Если $W_3$ ненулев, самосопряжён, калибровочно- и перестановочно-инвариантен
и имеет нулевые частичные следы, то каждый $K_{f_t}$ наследует эти свойства,
остаётся ненулевым при $k=0$ для $\lambda_3\ne0$ и является нелокальным
эффективным шеститочечным ядром.
\end{theorem}

\begin{nogocriterion}[Запрет объявлять форму выведенной]
Статическая нормировка, положительность спектральных весов и все внутренние
симметрии не выбирают форм-фактор: уже $f_1$ и $f_2$ различны, но проходят
один контракт. Поэтому нелокальный класс допущен, а конкретное ядро не
выведено. Следующий гейт обязан получить спектральную меру из одного
родителя либо сохранить $\lambda_3$, массы и вычеты внешними данными.
\end{nogocriterion}

\begin{protocol}
\begin{enumerate}
 \item поле вычислений: \textbf{$\mathbb Q(z,t)$};
 \item чисел с плавающей точкой: \textbf{$0$};
 \item точных представителей: \textbf{$2$};
 \item положительных спектральных атомов: \textbf{$1$ и $2$};
 \item общий статический предел: \textbf{$f(0)=1$};
 \item формы вне нуля совпадают: \textbf{нет};
 \item непрерывное семейство допустимых форм: \textbf{да};
 \item нелокальный класс: \textbf{допущен};
 \item конкретная спектральная мера: \textbf{не выведена};
 \item физический барионный полюс: \textbf{не выведен};
 \item следующий гейт: \textbf{родительское происхождение спектральной меры}.
\end{enumerate}
\end{protocol}

Машинный сертификат сохранён в
\path{s2t_v8_baryon_nonlocal_six_point_kernel_admission_gate_results.json}.